\documentclass[10pt,a4paper]{amsart}
\usepackage[margin=19mm]{geometry}
\usepackage{amsmath,amssymb,mathtools}
\usepackage[scr=rsfs,cal=euler]{mathalfa}
\usepackage{xfrac}
\usepackage[unicode,hidelinks,bookmarks=false]{hyperref}
\newcommand{\ie}{i.e.\ }
\newcommand{\cf}{cf.\ }
\newcommand{\resp}{respectively\ }
\newcommand{\Subsection}{\subsection}
\providecommand{\nobreakdash}{\nobreak\hbox{-}}
\setlength{\emergencystretch}{2em}
\multlinegap0pt
\newcommand{\Rpoly}[1]{\mathcal R_{#1}^{\vee}}
\theoremstyle{plain}
\newtheorem{theorem}{Theorem}[section]
\newtheorem{lemma}[theorem]{Lemma}
\newtheorem{proposition}[theorem]{Proposition}
\newtheorem{corollary}[theorem]{Corollary}
\theoremstyle{definition}
\newtheorem{definition}[theorem]{Definition}
\theoremstyle{remark}
\newtheorem{remark}[theorem]{Remark}
\theoremstyle{definition}
\newtheorem{problem}[theorem]{Problem}
\newtheorem{conjecture}[theorem]{Conjecture}
\title[A proof of the Athanasiadis-Chapoton ordinal-sum conjecture]{A proof of the Athanasiadis-Chapoton ordinal-sum conjecture}
\author{Congyi Luo}
\date{Source: arXiv:2609.28507v1 (CC BY 4.0)}
\begin{document}
\maketitle

\noindent\emph{Source and licence.} The delivered work is the article shown in the title above, version arXiv:2609.28507v1 (CC BY 4.0), distributed under the Creative Commons Attribution 4.0 International licence, as recorded in the arXiv licence metadata for this exact version and shown on the abstract page saved with this item. The full licence notice, which carries the licence URL and the address of that abstract page, is the bundled file \texttt{licenses/arxiv-2609.28507-CC-BY-4.0.txt}.
\par\smallskip

\noindent\textit{Editorial scope.} Counted unit: the article's Proposition (source label gauge) (source0.tex lines 209--219, source label \texttt{gauge}): ``For every nonempty finite preorder and every , (z) z z-z 1. In particular, for every real , m z: (z) m , ( Z E) Z 0 .''. It is delivered together with the article's complete original proof (source0.tex lines 220--258, one proof environment adjacent to the statement, 39 lines), copied line for line from the article's own text. Required context retained uncounted: envelope (lines 174--179); levels (lines 190--197). Editorial formatting only: an AMS \texttt{amsart} preamble that restates the article's own macro definitions and its own theorem declarations, and the theorem environments the retained lines use; the article's own class file is neither shipped nor input; the retained environments are renumbered sequentially in order of appearance, whereas the article prints them with its own numbering; the article's equation numbering is not reproduced and every cross-reference inside the retained text resolves to the wrapper's numbering. No citation occurs inside any retained line, so the wrapper carries no bibliography; All retained bodies are exact copies of the bundled \texttt{sources/211594/source0.tex}, so no omission receipt applies. No asset-inclusion command occurs inside any retained span and the package ships no image asset.


\section*{Retained context: envelope (source0.tex lines 174--179)}

\begin{lemma}\label{envelope}
The vector $\widehat z$ is the coordinatewise least among all nonnegative order-reversing vectors $v$ satisfying $v\ge z$, and
\begin{equation}\label{maxenvelope}
 \|\widehat z\|_\infty=\|z_+\|_\infty.
\end{equation}
\end{lemma}


\section*{Retained context: levels (source0.tex lines 190--197)}

\begin{lemma}\label{levels}
Let $v$ be a nonnegative order-reversing vector whose distinct positive coordinate values are $a_1<\cdots<a_s$. Put $a_0=0$ and $I_k=\{i:v_i\ge a_k\}$. Then each $I_k$ is a nonempty order ideal, and
\begin{equation}\label{leveldecomposition}
 v=\sum_{k=1}^s(a_k-a_{k-1})\mathbf1_{I_k},
 \qquad \sum_{k=1}^s(a_k-a_{k-1})=\|v\|_\infty.
\end{equation}
When $v=0$, both sums are empty.
\end{lemma}


\section{The counted unit: Proposition (source label gauge) and its complete original proof}

\begin{proposition}\label{gauge}
For every nonempty finite preorder $\tau$ and every $z\in\mathbb R^E$,
\begin{equation}\label{gdef}
 \gamma_\tau(z)=\|z_+\|_\infty+\|\widehat z-z\|_1.
\end{equation}
In particular, for every real $m\ge0$,
\begin{equation}\label{sublevel}
 m\Rpoly{\tau}=\{z:\gamma_\tau(z)\le m\},
 \qquad \gamma_\tau(\mathbb Z^E)\subseteq\mathbb Z_{\ge0}.
\end{equation}
\end{proposition}

\begin{proof}
Denote the right-hand side of~\eqref{gdef} by $\rho(z)$. Since $\widehat z\ge z$, the second term is $\sum_i(\widehat z_i-z_i)$, with each summand nonnegative. We first prove that $\Rpoly{\tau}=\{z:\rho(z)\le1\}$.

If $z\in\Rpoly{\tau}$, the convex-hull definition gives an expression
\[
 z=p-\beta,\qquad
 p=\sum_{\varnothing\ne I\in\mathcal J(\tau)}\alpha_I\mathbf1_I,
 \qquad \alpha_I,\beta_i\ge0,\quad
 \sum_I\alpha_I+\sum_i\beta_i=1,
\]
where $\beta=(\beta_i)_{i\in E}$. Put $\lambda=\sum_I\alpha_I$. The vector $p$ is nonnegative, order-reversing, and dominates $z$, so Lemma~\ref{envelope} gives $\widehat z\le p$. Also, $z_i\le p_i\le\lambda$ and $\lambda\ge0$, whence $\|z_+\|_\infty\le\lambda$. Therefore
\[
 \rho(z)\le\lambda+\sum_i(p_i-z_i)
          =\lambda+\sum_i\beta_i=1.
\]

Conversely, if $\rho(z)\le1$, Lemma~\ref{levels} and~\eqref{maxenvelope} give an expression
\[
 \widehat z=\sum_{k=1}^s\alpha_k\mathbf1_{I_k},
 \qquad \alpha_k\ge0,\qquad
 \sum_k\alpha_k=\|z_+\|_\infty.
\]
Consequently,
\[
 z=\sum_{k=1}^s\alpha_k\mathbf1_{I_k}
       +\sum_{i\in E}(\widehat z_i-z_i)(-e_i).
\]
All coefficients are nonnegative and sum to $\rho(z)\le1$. Adding the origin with coefficient $1-\rho(z)$ gives a convex combination in $\Rpoly{\tau}$, so $z\in\Rpoly{\tau}$. If $\widehat z=0$, the first sum is empty and the argument is unchanged.

For $c\ge0$, the definition gives $\widehat{cz}=c\widehat z$, hence $\rho(cz)=c\rho(z)$. Thus, for $m>0$,
\[
 z\in m\Rpoly{\tau}
 \ \Longleftrightarrow\ z/m\in\Rpoly{\tau}
 \ \Longleftrightarrow\ \rho(z)\le m.
\]
If $\rho(z)=0$, then $z_+=0$ and $\widehat z=z$. The first equality implies $z\le0$, while the second, together with $\widehat z\ge0$, gives $z=0$. The same sublevel-set identity therefore holds when $m=0$.

It follows that the set of scales $r$ for which $z\in r\Rpoly{\tau}$ is exactly $[\rho(z),\infty)$. Taking the infimum yields $\gamma_\tau(z)=\rho(z)$ and establishes~\eqref{sublevel}. Finally, for an integer input, $\widehat z$ is an integer vector and both terms in~\eqref{gdef} are nonnegative integers.
\end{proof}

\end{document}
